\section{Verification purpose}

This case verifies sensitivity of firebrand-driven WUI spread to structural spatial representation. Five simulation conditions cover the baseline grid, native fine and coarse grids, and aligned and misaligned integrated-structure representations under one physical geometry and source definition.

The test quantifies native cell-scale distortion and determines whether integrated representation restores structure-level accumulation, ignition chronology, and spread-rate consistency. Capability-gated simulation conditions remain explicit and are not replaced by a physically different native-cell calculation.

\section{Mathematical formulation and numerical implementation}

\subsection{Heat release and generation}

For a structure ignited at \(t_i\), the design heat-release-rate (HRR) per unit area (HRRPUA) is
\[
q''(\tau)=
\begin{cases}
q''_{\max}\tau/300,&0<\tau<300,\\
q''_{\max},&300\le\tau\le3900,\\
q''_{\max}(4200-\tau)/300,&3900<\tau<4200,\\
0,&\text{otherwise},
\end{cases}
\quad \tau=t-t_i,\qquad q''_{\max}=\SI{400}{kW.m^{-2}}.
\]
For footprint \(A_s\), \(\mathrm{HRR}_i=q''A_s\), and the emission rate is
\[
\dot N_i(t)=G_R'\mathrm{HRR}_i(t),\qquad
G_R'=\SI{10}{pcs.MW^{-1}.s^{-1}}.
\]
The configuration therefore uses per-MW generation; per-area generation would
not preserve heat-release-scaled production for this design fire.

\subsection{Himoto transport and analytical load}

The downwind distance is lognormal, with physical-space moments
\[
\mu_x/D=0.47(B^*)^{2/3},\qquad \sigma_x/D=0.88(B^*)^{1/3},
\]
\[
B^*=\frac{U}{\sqrt{gD}}
\left(\frac{\rho_p}{\rho_\infty}\right)^{-3/4}
\left(\frac{d_p}{D}\right)^{-3/4}
\left(\frac{\mathrm{HRR}}
{\rho_\infty c_pT_\infty g^{1/2}D^{5/2}}\right)^{1/2}.
\]
Here \(D=\SI{10}{m}\), \(\rho_p=\SI{100}{kg.m^{-3}}\), and
\(d_p=\SI{5}{mm}\), matching the reference formulation and current source. The postprocessor
independently evaluates the analytical deposited-load convolution
\[
N_k=\sum_{i<k}\int_0^T\dot N_i(t)
\left[\int_{x_k}^{x_{k+1}}
f(x-x_i;\mu_i(t),\sigma_i(t))\,dx\right]dt
\]
with a \SI{5}{s} quadrature interval. Deposited counts are divided by physical
footprint area, not a resolution-dependent raster-cell area.

\subsection{Physical ignition and expected chronology}

The configured ignition model uses
\(N_{\rm crit}=\SI{2928.9}{pcs.m^{-2}}\)~\cite{Qin2025}. At the structure
spanning \(20\le x\le\SI{30}{m}\), the historical analytical load reaches
this threshold near \SI{2496.9}{s}. The case-specific postprocessor preserves the
published composite full-ignition reference of \SI{2858.8}{s}; its difference
from direct addition of the currently configured delays is quantified
explicitly below.
Targets beyond the roughly \SI{80}{m} peak-heat-release rate (HRR) spotting range require
successive structure ignition and consequently have nonlinear arrival times.

\subsection{Resolution signatures}

On the \SI{2.5}{m} native mesh (native fine-grid simulation condition), one physical structure occupies
multiple independent cells. Uneven deposition and cell-local ignition should
produce a nonzero within-structure load range and mean rate of spread (ROS) near
\SI{0.001}{m.s^{-1}}, compared with about \SI{0.006}{m.s^{-1}} for the
exact-grid baseline.
On the \SI{30}{m} native mesh (native coarse-grid simulation condition), mode resampling inflates footprint
and HRR, causing artificially early spread; the documented response gives a
baseline-to-coarse load ratio near 1.86~\cite{Qin2025}.

The integrated-structure treatment groups all cells belonging to one structure,
\[
\bar N=\frac{\sum_{j\in s}N_j}{A_s},\qquad
\mathrm{HRR}_s=q''A_s,
\]
then redistributes load uniformly and emits from the structure centroid.
The integrated-aligned simulation condition expects an aligned \SI{2.5}{m} mesh to
reproduce the baseline. The integrated-misaligned simulation condition permits up to 33\%
load deviation and 10\% ignition-time deviation
on a misaligned \SI{3.24}{m} mesh.

\section{Derivation of expected behavior}
The exact grid should reproduce the analytical structure-load and ignition chronology. Native fine cells should expose within-structure imbalance, native coarse cells should expose footprint inflation, and integrated simulation conditions should recover the physical-structure response independently of alignment.

% BEGIN EXPLICIT EXPECTED-RESULT DERIVATION
At the \(\SI{400}{kW\,m^{-2}}\) plateau, one physical
\(10\times10~\mathrm{m}\) structure has
\[
 HRR_{\max}=400(100)=40000~\mathrm{kW}=40~\mathrm{MW},
 \qquad \dot N_{\max}=10(40)=400~\mathrm{pcs\,s^{-1}},
\]
and its critical total count is
\(N_{\mathrm{crit}}=2928.9(100)=292890\) particles.
Using the postprocessor's adjusted wind
\(U=0.447(40)/0.87=20.5517~\mathrm{m\,s^{-1}}\) at peak HRR gives
\[
 B^*=7.37372,\qquad
 \bar x=0.47(7.37372)^{2/3}(10)=17.8056~\mathrm{m},
\]
\[
 s_x=0.88(7.37372)^{1/3}(10)=17.1282~\mathrm{m},\qquad
 (\mu,\sigma)=(2.55195,0.80939).
\]
These values are inserted in the cell-cumulative distribution function (CDF) convolution at each 5-s midpoint;
the HRR dependence is recomputed rather than freezing the peak values.

There is an explicit provenance discrepancy in the historical chronology.
The case-specific constants give
\[
 2496.9+42.1+300=2839.0~\mathrm{s},
\]
whereas the preserved historical composite reference used by postprocessing
is \(2858.8~\mathrm{s}\), 19.8 s later.  The latter is therefore treated as
a published numerical estimate, not as the arithmetic result of the three
displayed constants.  The 19.8-s difference is only
\(19.8/2858.8=0.00693\), well inside the predeclared 10-percent baseline
ignition tolerance, so this provenance mismatch cannot change that component's
decision.

The exact-grid ideal errors are zero for analytical load and the chosen
historical ignition reference.  Native-grid expected signatures are a
fine/baseline ROS ratio near \(0.001/0.006=0.1667\) and a coarse load ratio
near 1.86.  Integrated aligned references are the exact-grid structure loads
and times; for the misaligned grid, one unresolved boundary strip is
\(3.24/10=0.324\), motivating the explicit 0.33 load allowance.
% END EXPLICIT EXPECTED-RESULT DERIVATION

\section{Verification metrics and acceptance criteria}

The postprocessor calculates every metric explicitly. The exact-grid baseline
compares structure-integrated load with the analytical convolution defined
above and second-structure ignition with \SI{2858.8}{s}. The native fine and
coarse simulation conditions verify the documented cell-local supporting results rather than claiming
convergence. The integrated aligned and misaligned simulation conditions compare
physical-structure profiles with the exact-grid baseline. Missing output or
capability is reported as
\emph{NOT EVALUABLE}, never as an implicit pass.

\subsection{Justification of metric quantification}

The criteria intentionally separate reproduction of native-grid defects from
accuracy requirements for the baseline and integrated treatments.  The 10\%
baseline load and ignition bands accommodate rasterized deposition and
threshold-crossing time while remaining small relative to the resolution
effects documented for this configuration~\cite{Qin2025}.  A native fine-grid
load range of at least 5\% is set above roundoff and output quantization, and a
ROS ratio no greater than 0.50 requires the material factor-of-two slowdown
that identifies independent-cell behavior rather than a visually detectable
but negligible difference.  The coarse-grid load ratio is centered on the
reported 1.86 response; its 25\% band permits the severe footprint snapping
intrinsic to a \SI{30}{m} cell without accepting a reversal of the expected
inflation.

For the aligned integrated grid, exact structure boundaries and conservation
of integrated load justify the stringent 1\% load limit.  Its 5\% ignition
limit is wider because ignition is a nonlinear threshold event recorded at a
discrete time.  For the \SI{3.24}{m} misaligned grid, one boundary-cell width
relative to a \SI{10}{m} structure is \(3.24/10=0.324\); the 0.33 maximum
load-error limit therefore admits approximately one unresolved boundary strip
but not a larger structure-scale bias.  The corresponding 10\% ignition bound
allows threshold amplification of that overlap error.  Capability and
completeness are conjunctive because the integrated claims cannot be inferred
from native-grid substitutes.

\subsection{Predeclared acceptance criteria}

Table~\ref{tab:acceptance-contract} summarizes the metrics fixed before
inspection of the calculated results. Numerical values and component statuses
are reported separately in Section~7.

\begin{table}[H]
\centering
\footnotesize
\caption{Predeclared verification metrics and acceptance criteria.}
\label{tab:acceptance-contract}
\begin{tabular}{@{}p{0.23\linewidth}p{0.47\linewidth}p{0.21\linewidth}@{}}
\toprule
Metric & Output, region, and aggregation & Acceptance criterion\\
\midrule
Baseline load error & Mean structure-load relative error against the analytical deposition reference, excluding the source structure. & \(\le0.10\) \\
Baseline ignition error & Relative error of second-structure ignition time against \SI{2858.8}{s}. & \(\le0.10\) \\
Native fine-grid load distortion & Median within-structure relative load range. & \(\ge0.05\) \\
Native fine-grid ROS ratio & Fine-grid fitted mean ROS divided by baseline ROS. & \(\le0.50\) \\
Native coarse-grid load ratio & Baseline structure-4 load divided by coarse structure-2 load. & \(1.86\pm25\%\) \\
Integrated aligned load error & Mean matched structure-load relative error against baseline. & \(\le0.01\) \\
Integrated aligned ignition error & Mean matched ignition-time relative error against baseline. & \(\le0.05\) \\
Integrated misaligned load error & Maximum matched structure-load relative error against baseline. & \(\le0.33\) \\
Integrated misaligned ignition error & Maximum matched ignition-time relative error against baseline. & \(\le0.10\) \\
Overall decision & All configured simulation conditions and metrics must be evaluable; unavailable integrated capability prevents a PASS. & PASS only if all rows pass \\
\bottomrule
\end{tabular}
\end{table}

\section{Simulation configuration and predicted outcome}

Figure~\ref{fig:input-configuration} records the prepared spatial inputs for the 10 m baseline grid.
These panels show the prepared spatial fields, remove the
two-cell numerical buffer, and retain map coordinates in metres.  The left panel shows
the most informative fuel or structure layout available for the case; the right panel
shows the cells selected by the non-positive initial level-set field, making a localized
ignition or firebrand-emission source explicit.

\begin{figure}[H]
\centering
\EvidenceFigure[width=0.98\textwidth,height=0.42\textheight,keepaspectratio]{../figures/input_configuration.pdf}
\caption{Whole-domain prepared input configuration for the representative simulation condition.  The
physical domain is shown without the numerical halo; coordinates, configured wind values,
fuel or structure layout, and initial ignition/source geometry are identified in the panels.}
\label{fig:input-configuration}
\end{figure}

\begin{table}[H]
\centering\small
\caption{Five simulation conditions generated from the common physical problem.}
\begin{tabular}{llllp{0.31\linewidth}}
\toprule
Simulation condition & \(\Delta x\) (m) & \(t_{\max}\) (s) & Treatment & Expected result\\
\midrule
Baseline & 10.00 & 22000 & native & analytical deposition and ignition\\
Native fine & 2.50 & 50000 & native fine & imbalance and slower ROS\\
Native coarse & 30.00 & 22000 & native coarse & inflated HRR and earlier spread\\
Integrated aligned & 2.50 & 22000 & integrated aligned & baseline-equivalent profile\\
Integrated misaligned & 3.24 & 22000 & integrated misaligned & bounded profile error\\
\bottomrule
\end{tabular}
\end{table}

Each nominal timestep is \(\Delta t_{\mathrm{nominal}}=0.5\Delta x/U\); the generated step is no larger than this bound. All structure rasters reference
a case-specific model 14 containing the design fire and physical delay. Structures
form a periodic \(220\times100\)-m two-dimensional lattice. Each represented
cell uses its true square area, so heat-release rate per unit area (HRRPUA) remains \SI{400}{kW.m^{-2}} and the
dimensionless ignition-hardening factor remains one at every resolution. Every
structure in the first upwind column is initialized at \(t=0\).
The global vegetative surface-spread setting switch remains true to suppress vegetative surface-fire propagation.  In the current source, initialized model-14 buildings independently refresh their transient HRR and therefore remain active under this setting.  The generated non-structure fuel model 93 is also nonburnable, providing a second explicit isolation of the building-only experiment.
Because the wind Courant--Friedrichs--Lewy (CFL) bound requires a subsecond step, preprocessing computes the
nominal fractional value.  When that value does not divide the requested duration,
it uses \(N=\lceil t_{\mathrm{requested}}/\Delta t_{\mathrm{nominal}}\rceil\)
and reduces the generated step to
\(\Delta t=t_{\mathrm{requested}}/N\).  The stop therefore remains the exact
requested duration and an exact timestep multiple; the simulation specification records nominal
and actual steps, CFL values, and the integer step count.

\subsection{Preprocessing, execution, and postprocessing}

Exact cell--footprint overlap defines the structure fractions and identities;
mode classification defines the native fine- and coarse-grid experiments.
Only model-supported structure representations are simulated. Final firebrand
counts are normalized by physical structure area and aggregated by structure
identity, together with the fire arrival times. Evaluation requires agreement
between the prescribed and simulated inputs, successful termination, unique
stop-time, timestep, and meteorological-interval values, and exact divisibility
of duration by timestep. One final output-time record and one corresponding
deposition/arrival-time pair are required. An output beyond the nominal stop
is accepted only if subtracting the meteorological interval gives a time on
the timestep grid within one timestep of the stop. Ambiguous, early, off-grid,
or duplicate evidence remains NOT EVALUABLE.

\clearpage
\section{Actual simulation results}

Figure~\ref{fig:domain-result} shows the representative time-of-arrival field from the 10 m baseline grid.
The displayed field is taken from the simulated results, masks missing values, and excludes
the two-cell numerical buffer.  This completed baseline provides whole-domain evidence even though simulation conditions requiring the unavailable integrated-structure capability remain not evaluable.

\begin{figure}[H]
\centering
\EvidenceFigure[width=0.96\textwidth,height=0.50\textheight,keepaspectratio]{../figures/domain_result.pdf}
\caption{Whole-domain time-of-arrival field from the representative completed ELMFIRE run.  This
domain-scale result supplements, rather than replaces, the higher-level verification
profiles, convergence plots, ensemble statistics, and calculated acceptance metrics below.}
\label{fig:domain-result}
\end{figure}

\CompletedVariantCount{} of \RequestedVariantCount{} simulation conditions currently have
complete, configuration-matched ELMFIRE output. Source-aligned physical-state
figures are provided for all three native representations; a figure explicitly
states when its current matching output is unavailable. The two
integrated-structure simulation conditions remain capability gated.

\begin{figure}[H]
\centering
\EvidenceFigure[width=\linewidth,height=0.68\textheight,keepaspectratio]{baseline_reference.pdf}
\caption{Exact-grid baseline observables: structure load, ignition chronology,
leading-edge trajectory, and instantaneous ROS. Analytical load and
second-structure timing references are overlaid where available.}
\label{fig:wsr-baseline}
\end{figure}

\begin{figure}[H]
\centering
\EvidenceFigure[width=\linewidth,height=0.68\textheight,keepaspectratio]{native_fine_reference.pdf}
\caption{Native fine-grid observables. Cell-scale load variation within each
physical structure is retained so discretization distortion is visible rather
than hidden by structure averaging.}
\label{fig:wsr-fine}
\end{figure}

\begin{figure}[H]
\centering
\EvidenceFigure[width=\linewidth,height=0.68\textheight,keepaspectratio]{native_coarse_reference.pdf}
\caption{Native coarse-grid observables. Mode resampling changes the represented
footprint and separation, and the resulting load and ignition histories are
shown on their actual physical coordinates.}
\label{fig:wsr-coarse}
\end{figure}

Figure~\ref{fig:wsr-summary} is the supplemental cross-simulation condition diagnostic. The
native-test decision is \textbf{\NativeStatus}; the complete-suite decision is
\textbf{\OverallStatus}.

\begin{figure}[H]
\centering
\EvidenceFigure[width=\linewidth]{wui_structure_resolution.pdf}
\caption{Structural-resolution summary. (a) Rasterized layouts; (b) deposited
load by physical structure; (c) ignition chronology; and (d) component decision
status, with unsupported components shown as not evaluated.}
\label{fig:wsr-summary}
\end{figure}

\section{Calculated metrics and verification decision}

\begin{table}[H]
\centering\small
\caption{Calculated verification metrics and component decisions.}
\begin{tabular}{@{}p{0.37\linewidth}p{0.18\linewidth}p{0.17\linewidth}p{0.14\linewidth}@{}}
\toprule
Metric & Acceptance limit & Calculated & Status\\
\midrule
\MetricRows
\bottomrule
\end{tabular}
\end{table}

The cited model description~\cite{Qin2025} states that the integrated-structure method
had not yet been implemented in ELMFIRE. The preprocessor confirms this against
the current source. The integrated aligned and misaligned cases are therefore
generated as prospective simulation conditions
but are not run while that feature is absent. Substituting the native
cell-local formulation would change the prescribed experiment and could
produce misleading conclusions about resolution independence. Earlier input
configurations are excluded from the present comparison.

\section{Technical interpretation}
The corrected exact-grid baseline and native coarse simulation condition both terminate at
an exact \SI{22000}{s} final dump with matching input identity records.  The baseline
load error is 0.5435 and its ignition-time error is 0.1861, both above the
0.10 limits; these are genuine numerical failures.  The coarse
baseline-to-coarse load ratio is 1.622 and passes the prescribed
\(1.86\pm25\%\) band.

The native fine-grid member was executed toward its prescribed
\SI{50000}{s} endpoint but could not complete safely.  It reached
\SI{36000.065}{s}; resident memory grew from approximately 204 MB near
\SI{15000}{s} to 960 MB, while successive \SI{500}{s} dump intervals slowed
from roughly 198 wall seconds to about 11 minutes.  The run was stopped before
exhausting the shared environment, and no completion input identity record was written,
so its partial rasters are correctly excluded.  This reproduces an
implementation resource-scaling limitation in the active ember-tracking path,
not a stop-time or raster-configuration defect.

The integrated simulation conditions test whether one physical structure remains the source
and collector; their missing capability is a source-interface limitation and
is likewise reported as not evaluable rather than as a numerical pass or
failure.  Consequently the native-test status is \textbf{FAIL}, while the full
suite remains \textbf{NOT EVALUABLE (capability missing)}.

\section{Scope and limitations}
The case is an idealized two-dimensional WUI lattice with fixed structure geometry and wind. A representative cross-stream lattice row supplies the one-dimensional structure profiles, while whole-domain maps expose unintended lateral behavior. The case verifies resolution behavior and capability coverage, not the physical accuracy of the structural transport or ignition correlations.

\begin{thebibliography}{9}
\bibitem{Qin2025}
Y. Qin, \emph{A Physics-Based Eulerian Framework for Modeling Firebrand
Showering in Regional-Scale Wildland and Wildland--Urban Interface Fire
Simulations}, Ph.D. dissertation, University of Maryland, College Park, 2025.
\end{thebibliography}
